%% Electronic Journal of Probability article starter.
\documentclass[EJP,preprint]{ejpecp}

\TITLE{A coupling argument for a finite-state sampler}
\SHORTTITLE{A coupling argument for a finite-state sampler}
\AUTHORS{Ada Lovelace and Grace Hopper}
\KEYWORDS{coupling; Markov chains; sampling}
\AMSSUBJ{60J10}
\ABSTRACT{We give a compact worked example of a coupling estimate for a
finite-state Markov chain. The document demonstrates the metadata and
theorem structure expected by the Electronic Journal of Probability while
remaining short enough to adapt as a preprint.}
\DOI{10.0000/example}
\ARXIVID{2601.00001}

\begin{document}

\begin{equation}
  \mathbb{P}(X_t\ne Y_t)\leq (1-\varepsilon)^t,
  \label{eq:coupling}
\end{equation}
for a pair of chains that share a successful update with probability at least
$\varepsilon$ at each step.

\section{A finite-state example}
Let $P$ be a transition matrix on a finite state space and suppose that two
copies of the chain use the same random update whenever they are in distinct
states. If one update coalesces them with probability at least
$\varepsilon$, the failure probability after $t$ updates is bounded by the
right-hand side of~\eqref{eq:coupling}.

\begin{theorem}
Under the uniform coalescence assumption, the coupled chains meet almost
surely and their total-variation distance is at most
$(1-\varepsilon)^t$ at time $t$.
\end{theorem}

\begin{proof}
The probability that no coalescing update has occurred in the first $t$
steps is at most $(1-\varepsilon)^t$. After the first meeting, use identical
updates, so the chains remain together. The coupling inequality gives the
claimed total-variation bound.
\end{proof}

\section{Practical checklist}
For an applied manuscript, record the state space, update rule, coupling
event, and the value of $\varepsilon$. Separate those assumptions from the
numerical results so that a reader can reproduce the argument.

\begin{table}[ht]
  \centering
  \caption{Illustrative convergence diagnostics.}
  \begin{tabular}{lrr}
    \hline
    Steps $t$ & Bound & Measured distance \\
    \hline
    4  & 0.41 & 0.36 \\
    12 & 0.08 & 0.06 \\
    24 & 0.01 & 0.01 \\
    \hline
  \end{tabular}
\end{table}

\section{Conclusion}
Replace this worked example with the statement, proof, and computational
details of your own probability article.

\begin{acks}
The authors thank the probability seminar for a careful reading of the
coupling argument.
\end{acks}

\begin{thebibliography}{9}
\bibitem{lindvall2002lectures}
Lindvall, T. (2002). \emph{Lectures on the Coupling Method}. Dover.
\end{thebibliography}

\end{document}
